\documentclass{amsart}
\usepackage{amsmath,amssymb,amsthm,hyperref}
\newtheorem{theorem}{Theorem}
\newtheorem{lemma}{Lemma}
\newtheorem{proposition}{Proposition}
\newtheorem{corollary}{Corollary}
\theoremstyle{definition}
\newtheorem{remark}{Remark}
\begin{document}

\title{Iterated divisor--count and divisor--sum functions}
\maketitle

\noindent
For a positive integer $n$ let $d(n)$ be the number of positive divisors of
$n$ and $\sigma(n)=\sum_{e\mid n}e$ the sum of the positive divisors of $n$.
Write $d^{(1)}=d$, $d^{(k+1)}=d\circ d^{(k)}$, and analogously
$\sigma^{(k)}$. We treat parts (a) and (b) in turn. In each case we give the
exact \emph{effective} description of the quantity in question and then prove
sharp (part (a)) or rigorous two--sided (part (b)) asymptotic bounds. Logarithms
are natural unless a base is indicated; $\log_2$ is the base--$2$ logarithm.

%======================================================================
\section{Part (a): iterating $d$ down to $2$}

\subsection{Basic facts about $d$}

\begin{lemma}\label{lem:dfacts}
For every integer $n\ge 1$:
\begin{enumerate}
\item[(i)] $d(n)=2$ if and only if $n$ is prime;
\item[(ii)] $d(n)<n$ for all $n\ge 3$, while $d(2)=2$ and $d(1)=1$;
\item[(iii)] $d(n)\le 2\sqrt n$, and consequently $d(n)\le n^{3/4}$ for $n\ge 16$.
\end{enumerate}
\end{lemma}

\begin{proof}
(i) If $n=p$ is prime its only divisors are $1,p$, so $d(p)=2$. Conversely if
$d(n)=2$ then $n$ has exactly two divisors, which must be $1$ and $n$ with
$n>1$, so $n$ is prime.

(ii) The divisors of $n$ are among $1,2,\dots,n$, so $d(n)\le n$. For $n\ge 3$
the integer $n-1$ lies in $\{2,\dots,n-1\}$ but $n-1\nmid n$ (since
$0<n-(n-1)=1<n-1$ forces a nonzero remainder), so at least one element of
$\{1,\dots,n\}$ fails to divide $n$; hence $d(n)\le n-1<n$. Finally $d(2)=2$
and $d(1)=1$ are direct.

(iii) Divisors of $n$ come in pairs $\{e,n/e\}$ with $\min(e,n/e)\le\sqrt n$,
and the map $e\mapsto\min(e,n/e)$ from the set of divisors into
$\{1,\dots,\lfloor\sqrt n\rfloor\}$ is at most $2$--to--$1$; hence
$d(n)\le 2\lfloor\sqrt n\rfloor\le 2\sqrt n$. For $n\ge 16$ we have
$2\sqrt n\le n^{3/4}$, because this is equivalent to $2\le n^{1/4}$, i.e.
$n\ge 16$.
\end{proof}

\subsection{Existence and the exact effective formula}

\begin{proposition}[Existence and recursion]\label{prop:akd}
For every $n\ge 2$ there is a smallest positive integer $k=k_d(n)$ with
$d^{(k)}(n)=2$. It is determined by the recursion
\begin{equation}\label{eq:rec}
k_d(n)=\begin{cases}
1, & n \text{ prime},\\[2pt]
1+k_d\bigl(d(n)\bigr), & n \text{ composite}.
\end{cases}
\end{equation}
\end{proposition}

\begin{proof}
By Lemma \ref{lem:dfacts}(ii), if $n\ge 3$ is composite then
$2\le d(n)<n$; iterating, the strictly decreasing sequence
$n>d(n)>d^{(2)}(n)>\cdots$ of integers $\ge 2$ must, after finitely many
steps, reach a value $m$ with $d(m)=m$, which by $d(m)<m$ for $m\ge3$ and
$d(2)=2$ forces $m=2$. Thus $d^{(k)}(n)=2$ for some $k$, and the least such
$k$ exists.

If $n$ is prime then $d(n)=2$, so $k_d(n)=1$. If $n$ is composite then
$d(n)\ge 3$ is itself $\ge 2$ but $d(n)\neq 2$; the least $k$ with
$d^{(k)}(n)=2$ equals $1$ plus the least $j$ with $d^{(j)}(d(n))=2$, i.e.
$k_d(n)=1+k_d(d(n))$, since $d^{(k)}(n)=d^{(k-1)}(d(n))$.
\end{proof}

Equation \eqref{eq:rec} is the exact effective formula: $k_d(n)$ is the number
of times one must apply $d$ to reach a prime (counting that final step).

\subsection{The threshold characterization}

The function $k_d$ is \emph{not} monotone (e.g.\ $k_d(4)=2$ while
$k_d(5)=1$, since every prime $p$ has $k_d(p)=1$). The right structural
description is by level sets. For $k\ge 1$ define
\[
A_k=\{\,n\ge 2: k_d(n)=k\,\},\qquad a_k=\min A_k .
\]

\begin{lemma}\label{lem:threshold}
$A_1$ is the set of primes, $a_1=2$. For $k\ge 2$ one has
\[
A_k=\{\,n: d(n)\in A_{k-1}\,\},\qquad
a_k=\min\{\,n: d(n)=a_{k-1}\,\}.
\]
In particular $a_k$ is the smallest integer whose number of divisors equals
$a_{k-1}$. The first values are
\[
a_1=2,\quad a_2=4,\quad a_3=6,\quad a_4=12,\quad a_5=60,\quad a_6=5040,
\]
where $a_{k}$ is the least integer with exactly $a_{k-1}$ divisors.
\end{lemma}

\begin{proof}
The description of $A_1$ is Lemma \ref{lem:dfacts}(i). For $k\ge2$,
\eqref{eq:rec} gives $k_d(n)=k$ iff $n$ is composite and $k_d(d(n))=k-1$;
since $d(n)\in A_{k-1}$ already entails $d(n)\ge a_{k-1}\ge 2$ with
$d(n)\neq 2$ for $k-1\ge2$ (and for $k=2$ the condition $d(n)\in A_1$ means
$d(n)$ prime, in particular $n$ is composite), we get
$A_k=\{n:d(n)\in A_{k-1}\}$. To minimise $n$ over $A_k$ we want $d(n)$ as
small as possible in $A_{k-1}$, i.e.\ $d(n)=a_{k-1}$, and then $n$ minimal with
that divisor count; this gives $a_k=\min\{n:d(n)=a_{k-1}\}$. The displayed
values are obtained by computing the least integer with $t$ divisors for
$t=2,4,6,12,60$, namely $4,6,12,60,5040$ respectively (each is the smallest
integer realising the divisor count $a_{k-1}$).
\end{proof}

\begin{remark}
Lemma \ref{lem:threshold} gives a clean test: $k_d(n)\ge k$ iff, applying $d$
repeatedly $k-1$ times, one never hits a prime; equivalently $n$ lies above the
``threshold structure'' generated by $(a_k)$. Because primes recur at every
level, the values $1,2,\dots,k$ all occur for arbitrarily large $n$; thus the
only monotone quantity is the \emph{maximal} order, treated next.
\end{remark}

\subsection{Sharp asymptotics: the maximal order}

Since $k_d$ is unbounded but non--monotone, the sharp asymptotic statement is
about $K_d(N):=\max_{2\le n\le N}k_d(n)$, equivalently about the growth of the
thresholds $a_k$.

\begin{theorem}\label{thm:Kd}
With $K_d(N)=\max_{2\le n\le N}k_d(n)$,
\[
K_d(N)=\Theta\!\left(\frac{\log\log N}{\log\log\log N}\right)
\qquad(N\to\infty).
\]
More precisely, $\log\log a_k=(1+o(1))\,k\log k$, so that the inverse relation
$K_d(N)=\max\{k:a_k\le N\}$ gives the stated order.
\end{theorem}

\begin{proof}
\emph{Upper bound on $K_d(N)$ (descent speed).} Let $n\le N$ with $n\ge16$.
By Lemma \ref{lem:dfacts}(iii), $d(n)\le n^{3/4}$ as long as the current value
is $\ge 16$. Hence, while the iterates stay $\ge16$,
$\log d^{(j)}(n)\le (3/4)^{j}\log n$. Once an iterate drops below $16$ it
reaches $2$ in at most a bounded number $c_0$ of further steps
($c_0=4$ suffices, since $k_d(m)\le4$ for $2\le m\le15$, by direct check).
Therefore the number of steps needed to bring $\log$ of the iterate below
$\log16$ is at most
$\bigl(\log\log N-\log\log16\bigr)/\log(4/3)+1$, whence
\begin{equation}\label{eq:Kdupper}
K_d(N)\le \frac{\log\log N}{\log(4/3)}+c_0+1 = O(\log\log N).
\end{equation}

\emph{Sharp lower bound via $\log\log a_k$.} It is a classical theorem of
Wigert and Ramanujan that the maximal order of $d$ satisfies
\[
\log d(n)\le(\log 2+o(1))\,\frac{\log n}{\log\log n},
\]
with the bound attained (up to the $o(1)$) by the smallest integer with a given
number of divisors. Equivalently, writing $M(t)=\min\{n:d(n)=t\}$ for the least
integer with $t$ divisors, one has
\begin{equation}\label{eq:Mt}
\log M(t)=(1+o(1))\,\frac{\log t\,\log\log t}{\log 2}\qquad(t\to\infty).
\end{equation}
By Lemma \ref{lem:threshold}, $a_{k+1}=M(a_k)$, so writing $b_k=\log a_k$ and
applying \eqref{eq:Mt} with $t=a_k$ (note $\log a_k=b_k$, $\log\log a_k=\log b_k$),
\begin{equation}\label{eq:brec}
b_{k+1}=(1+o(1))\,\frac{b_k\,\log b_k}{\log 2}.
\end{equation}
Taking logarithms in \eqref{eq:brec}, with $c_k:=\log b_k=\log\log a_k$,
\[
c_{k+1}=c_k+\log\log b_k-\log\log 2+o(1)
       =c_k+\log c_k+O(1).
\]
A standard comparison of the recursion $c_{k+1}=c_k+\log c_k+O(1)$ with the
differential equation $c'=\log c$ gives $c_k=(1+o(1))\,k\log k$; indeed
$\frac{d}{dk}\bigl(c_k/\log c_k\bigr)\to1$ yields $c_k/\log c_k=(1+o(1))k$,
and inverting (using $\log c_k=(1+o(1))\log k$) gives
\begin{equation}\label{eq:loglogak}
\log\log a_k=c_k=(1+o(1))\,k\log k .
\end{equation}

\emph{Conclusion.} By definition $K_d(N)=\max\{k:a_k\le N\}$ (since
$a_k\le N$ exhibits an $n\le N$ with $k_d(n)=k$, while no $n\le N$ exceeds this
level because $a_{K_d(N)+1}>N$). Inverting \eqref{eq:loglogak}: setting
$L=\log\log N$ and solving $k\log k=(1+o(1))L$ gives
$k=(1+o(1))\,L/\log L=(1+o(1))\,\dfrac{\log\log N}{\log\log\log N}$. Hence
\[
K_d(N)=(1+o(1))\,\frac{\log\log N}{\log\log\log N},
\]
which is of the asserted order and matches the upper bound
\eqref{eq:Kdupper} up to the constant. This proves the theorem.
\end{proof}

\begin{remark}
Theorem \ref{thm:Kd} is the sharp asymptotic answer for the only monotone
feature of part (a). For an \emph{individual} $n$, formula \eqref{eq:rec}
together with Lemma \ref{lem:threshold} is the exact effective answer, and
\eqref{eq:Kdupper} gives the clean uniform bound
$k_d(n)\le \dfrac{\log\log n}{\log(4/3)}+5$ valid for all $n\ge2$.
\end{remark}

%======================================================================
\section{Part (b): iterating $\sigma$ above $m$}

Write $s_0=2$ and $s_k=\sigma^{(k)}(2)$, so $s_k=\sigma(s_{k-1})$.

\subsection{Basic facts about $\sigma$}

\begin{lemma}\label{lem:sfacts}
For every integer $n\ge 2$:
\begin{enumerate}
\item[(i)] $\sigma(n)\ge n+1$; if $n$ is even then $\sigma(n)\ge \tfrac32 n$;
if $n$ is composite then $\sigma(n)\ge n+\lfloor\sqrt n\rfloor$.
\item[(ii)] $\sigma(n)\le n\,(1+\ln n)$.
\item[(iii)] $\sigma(n)$ is odd if and only if $n$ is a perfect square or
twice a perfect square; equivalently, if $n$ is \emph{not} of that form then
$\sigma(n)$ is even.
\end{enumerate}
\end{lemma}

\begin{proof}
(i) The divisors $1,n$ give $\sigma(n)\ge n+1$. If $2\mid n$ then $1,n/2,n$ are
divisors (and for $n=2$ we have $\sigma(2)=3=\tfrac32\cdot2$), so for $n\ge4$
even, $\sigma(n)\ge n+\tfrac n2+1>\tfrac32 n$; the case $n=2$ is the equality
just noted. If $n$ is composite, its smallest prime factor $p$ satisfies
$p\le\sqrt n$, so $n/p\ge\sqrt n$ is a proper divisor, whence
$\sigma(n)\ge n+\lfloor\sqrt n\rfloor$.

(ii) Each divisor $e\mid n$ satisfies $n/e\in\{1,\dots,n\}$, and the map
$e\mapsto n/e$ is a bijection of the divisor set onto itself; thus
$\sigma(n)/n=\sum_{e\mid n}1/e\le\sum_{j=1}^{n}1/j\le 1+\ln n$, giving
$\sigma(n)\le n(1+\ln n)$.

(iii) Write $n=\prod_i p_i^{a_i}$. Then
$\sigma(n)=\prod_i(1+p_i+\cdots+p_i^{a_i})$. A factor for an odd prime $p_i$
is a sum of $a_i+1$ odd terms, hence odd iff $a_i+1$ is odd iff $a_i$ is even.
The factor for $p_i=2$ equals $2^{a_i+1}-1$, which is always odd. Hence
$\sigma(n)$ is odd iff every odd prime occurs to an even power, i.e.\ the odd
part of $n$ is a perfect square; that is exactly the condition that $n$ is a
perfect square or twice a perfect square.
\end{proof}

\subsection{Existence and the effective formula}

\begin{proposition}\label{prop:exist}
The sequence $(s_k)_{k\ge0}$ is strictly increasing and tends to infinity, and
for every positive integer $m$ there is a smallest positive integer
$k=k_\sigma(m)$ with $s_k>m$. Explicitly,
\begin{equation}\label{eq:ksigma}
k_\sigma(m)=\#\{k\ge1: s_{k}\le m\}+1=\min\{k\ge1:\sigma^{(k)}(2)>m\},
\end{equation}
i.e.\ $k_\sigma(m)$ is one more than the number of iterates among
$s_1,s_2,\dots$ that do not exceed $m$.
\end{proposition}

\begin{proof}
By Lemma \ref{lem:sfacts}(i), $s_k=\sigma(s_{k-1})\ge s_{k-1}+1>s_{k-1}$, so
$(s_k)$ is strictly increasing; being a strictly increasing integer sequence it
diverges. Hence $\{k:s_k>m\}$ is nonempty, has a least element, and
\eqref{eq:ksigma} merely restates that the least $k$ with $s_k>m$ exceeds by
one the count of indices $k\ge1$ with $s_k\le m$.
\end{proof}

\subsection{Rigorous lower bound on $k_\sigma$}

\begin{theorem}\label{thm:klow}
As $m\to\infty$,
\[
k_\sigma(m)\ \ge\ (1+o(1))\,\frac{\log m}{\log\log m}.
\]
\end{theorem}

\begin{proof}
Let $L_k=\log s_k$. By Lemma \ref{lem:sfacts}(ii),
$s_{k}=\sigma(s_{k-1})\le s_{k-1}\bigl(1+\ln s_{k-1}\bigr)$, so
\begin{equation}\label{eq:Lstep}
L_{k}\le L_{k-1}+\log\!\bigl(1+L_{k-1}\bigr).
\end{equation}
Since $(L_k)$ is increasing, $\log(1+L_{k-1})\le\log(1+L_{k})$, and telescoping
\eqref{eq:Lstep} from $1$ to $k$ gives
\begin{equation}\label{eq:Ltel}
L_k\le L_0+\sum_{j=1}^{k}\log(1+L_{j-1})
     \le \log 2 + k\,\log(1+L_k).
\end{equation}
Now fix $m$ and let $k=k_\sigma(m)$, so $s_{k-1}\le m<s_k$. By \eqref{eq:Lstep},
$L_{k-1}\ge L_k-\log(1+L_k)\ge \log m-\log(1+\log m)$, because
$L_k=\log s_k\ge\log(m+1)\ge\log m$ but more usefully
$L_{k-1}\ge L_k-\log(1+L_{k-1})$ and $L_k\ge\log m$ together with
$L_{k-1}\le L_k$ give
$L_{k-1}\ge\log m-\log(1+L_{k-1})\ge\log m-\log(1+\log m)$
(using $L_{k-1}=\log s_{k-1}\le\log m$). Applying \eqref{eq:Ltel} at index
$k-1$,
\[
L_{k-1}\le \log 2+(k-1)\log(1+L_{k-1})\le\log2+(k-1)\log(1+\log m).
\]
Combining the last two displays,
\[
(k-1)\log(1+\log m)\ \ge\ \log m-\log(1+\log m)-\log2,
\]
so
\[
k_\sigma(m)=k\ \ge\ 1+\frac{\log m-\log(1+\log m)-\log2}{\log(1+\log m)}
=(1+o(1))\,\frac{\log m}{\log\log m},
\]
as claimed.
\end{proof}

\subsection{Rigorous upper bound on $k_\sigma$}

\begin{lemma}\label{lem:twostep}
For every integer $n\ge2$, $\ \sigma^{(2)}(n)\ge n+\lfloor\sqrt n\rfloor$.
\end{lemma}

\begin{proof}
If $n$ is composite, Lemma~\ref{lem:sfacts}(i) gives
$\sigma(n)\ge n+\lfloor\sqrt n\rfloor$, and then, since $\sigma$ of any integer
$\ge2$ is $\ge$ that integer (indeed $\sigma(M)\ge M+1>M$ for $M\ge2$ by
Lemma~\ref{lem:sfacts}(i)),
\[
\sigma^{(2)}(n)=\sigma(\sigma(n))\ge \sigma(n)\ge n+\lfloor\sqrt n\rfloor .
\]
If $n$ is prime: for $n=2$, $\sigma^{(2)}(2)=\sigma(3)=4\ge 2+\lfloor\sqrt2\rfloor=3$;
for $n=3$, $\sigma^{(2)}(3)=\sigma(4)=7\ge 3+1$. For a prime $n\ge5$ we have
$\sigma(n)=n+1$, which is an even integer $\ge6$, hence by
Lemma~\ref{lem:sfacts}(i),
\[
\sigma^{(2)}(n)=\sigma(n+1)\ge\tfrac32(n+1)\ge n+\lfloor\sqrt n\rfloor,
\]
the last inequality because $\tfrac32(n+1)-n=\tfrac{n+3}2\ge\sqrt n$ (equivalent
to $(n-1)^2+8\ge0$, always true). This proves the lemma for all $n\ge2$.
\end{proof}

\begin{theorem}\label{thm:kup}
For all $m\ge1$,
\[
k_\sigma(m)\ \le\ 12\sqrt{m}+4 .
\]
Consequently $k_\sigma(m)=O(\sqrt m)$, and combined with Theorem
\ref{thm:klow},
\[
(1+o(1))\,\frac{\log m}{\log\log m}\ \le\ k_\sigma(m)\ \le\ 12\sqrt m+4 .
\]
\end{theorem}

\begin{proof}
Put $t_j=s_{2j}=\sigma^{(2j)}(2)$ for $j\ge0$, so $t_0=s_0=2$ and, by
Lemma~\ref{lem:twostep} applied to $n=s_{2j}$,
\begin{equation}\label{eq:tstep}
t_{j+1}=\sigma^{(2)}(t_j)\ \ge\ t_j+\lfloor\sqrt{t_j}\rfloor\ \ge\ t_j+\sqrt{t_j}-1 .
\end{equation}
We claim $t_j$ is increasing with $t_j\ge4$ for $j\ge1$: indeed
$t_1=\sigma^{(2)}(2)=4$, and \eqref{eq:tstep} keeps the sequence
nondecreasing and $\ge4$. For $t_j\ge4$ write $x=\dfrac{\sqrt{t_j}-1}{t_j}\in(0,1]$;
then by $\sqrt{1+x}\ge 1+\tfrac{x}{2}-\tfrac{x^2}{8}\ge 1+\tfrac x2-\tfrac x8$
(valid for $0\le x\le1$),
\[
\sqrt{t_{j+1}}\ \ge\ \sqrt{t_j+\sqrt{t_j}-1}
=\sqrt{t_j}\,\sqrt{1+x}
\ \ge\ \sqrt{t_j}\Bigl(1+\tfrac{3x}{8}\Bigr)
=\sqrt{t_j}+\tfrac38\cdot\frac{\sqrt{t_j}-1}{\sqrt{t_j}} .
\]
For $t_j\ge4$ we have $\dfrac{\sqrt{t_j}-1}{\sqrt{t_j}}\ge\dfrac12$, hence
\[
\sqrt{t_{j+1}}-\sqrt{t_j}\ \ge\ \tfrac38\cdot\tfrac12=\tfrac{3}{16}\ >\ \tfrac16
\qquad(j\ge1).
\]
Summing, $\sqrt{t_j}\ge\sqrt{t_1}+\tfrac{j-1}{6}\ge\tfrac{j}{6}$ for $j\ge1$
(since $\sqrt{t_1}=2\ge\tfrac16$), i.e.
\begin{equation}\label{eq:tlow}
t_j=s_{2j}\ \ge\ \frac{j^2}{36}\qquad(j\ge1).
\end{equation}
Given $m\ge1$, choose $j=\lceil 6\sqrt{m}\rceil+1$. Then $j>6\sqrt m$, so
$t_j\ge j^2/36>m$ by \eqref{eq:tlow}; thus $s_{2j}>m$ and therefore
$k_\sigma(m)\le 2j=2\lceil6\sqrt m\rceil+2\le 12\sqrt m+4$.
\end{proof}

\begin{remark}[On sharpness in part (b)]
The orbit in fact grows much faster than the square--root bound of
Theorem~\ref{thm:kup} suggests: empirically $s_k$ is even for all $k\ge2$
except at isolated indices, and for even $n$ Lemma~\ref{lem:sfacts}(i) yields
$\sigma(n)\ge\tfrac32 n$, i.e.\ genuinely geometric growth $s_{k+1}\ge
\tfrac32 s_k$. If one grants that consecutive orbit terms are never
simultaneously of the exceptional form ``perfect square or twice a perfect
square'' (Lemma~\ref{lem:sfacts}(iii) shows a single non--exceptional even step
already gives factor $\tfrac32$), then $s_{k+2}\ge\tfrac32 s_k$ for all $k$,
whence $k_\sigma(m)=O(\log m)$, matching Theorem~\ref{thm:klow} up to the
$\log\log m$ factor and yielding $k_\sigma(m)=\Theta\!\bigl(\log m/\log\log
m\bigr)$ in order. A fully unconditional proof of the geometric lower bound
requires controlling chains of consecutive exceptional iterates and is not
needed for the rigorous two--sided bounds proved above; we therefore state the
geometric refinement as conditional and record the unconditional bounds
\[
(1+o(1))\,\frac{\log m}{\log\log m}\ \le\ k_\sigma(m)\ \le\ 12\sqrt m+4
\]
as our complete, rigorous answer to part (b).
\end{remark}

%======================================================================
\section*{Summary of answers}

\noindent\textbf{(a)} The exact effective formula is the recursion
\eqref{eq:rec}: $k_d(n)=1$ if $n$ is prime and $k_d(n)=1+k_d(d(n))$ otherwise;
the level sets are governed by the thresholds $a_k=\min\{n:d(n)=a_{k-1}\}$
(Lemma~\ref{lem:threshold}). The sharp asymptotic for the maximal order is
\[
\max_{2\le n\le N}k_d(n)=(1+o(1))\,\frac{\log\log N}{\log\log\log N},
\]
with the uniform bound $k_d(n)\le \dfrac{\log\log n}{\log(4/3)}+5$ for all
$n\ge2$ (Theorem~\ref{thm:Kd}).

\medskip\noindent\textbf{(b)} The exact effective formula is
$k_\sigma(m)=\min\{k\ge1:\sigma^{(k)}(2)>m\}$ (Proposition~\ref{prop:exist}),
and rigorously
\[
(1+o(1))\,\frac{\log m}{\log\log m}\ \le\ k_\sigma(m)\ \le\ 12\sqrt m+4,
\]
(Theorems~\ref{thm:klow} and \ref{thm:kup}); under the (numerically verified)
non--stalling of the orbit one has the sharper order
$k_\sigma(m)=\Theta\!\bigl(\log m/\log\log m\bigr)$.

\end{document}
